% GENERATED FILE. Edit canonical lessons, then run npm run generate:books.
% canonical-source-hash: fnv1a64:447821bc2b2bcd4e
% canonical-lessons: ZH-C52-shou, ZH-W52-shou, ZH-C52-tou, ZH-C52-jiao, ZH-C52-yanjing, ZH-C52-erduo, ZH-R52-review-48-52

\chapter{Head, Hands and Feet}
\label{ch:zh-head-hands-and-feet}
\hlchaptermodality{52}

\begin{chapteropening}
\textbf{By the end of this chapter:} \emph{I can name my head, hands, feet, eyes and ears.}

Complete the last lesson of chapter 52: I can name my head, hands, feet, eyes and ears.
\end{chapteropening}

\section[\texorpdfstring{\zh{手} (shǒu)}{手 (shǒu)}]{\zh{手} (shǒu) — hand}

\label{lesson:ZH-C52-shou}



\begin{quote}
Before the new one: say the Mandarin for hot, then the Mandarin for to cry.
\end{quote}

\subsection*{You'll want to know: \zh{手}}
\textbf{\zh{手}} — \emph{shǒu} — \textquotedblleft{}hand\textquotedblright{}.

\begin{quote}
\textbf{\zh{我的手很小}} \emph{wǒ de shǒu hěn xiǎo} — My hands are small.
\end{quote}

The new character, in pieces: \zh{手} — a short fall, two bars, and a vertical that hooks at the foot. The next lesson puts it in your hand.

\begin{cousinweb}
\textbf{\zh{手}} (\emph{shǒu}) is inherited from Middle Chinese \emph{syuwX} (Guangyun, in Baxter's transcription), glossed with the pair 'hand and foot'.
\end{cousinweb}

\begin{grammarlens}[title={What you've built}]
One more word for saying how you are.
\end{grammarlens}

\subsection*{Your turn}
\begin{itemize}
\raggedright
  \item \emph{Say it:} \emph{shǒu}
  \item \emph{Say it:} \emph{shǒu}, once more
  \item \emph{Say it:} \emph{wǒ de shǒu hěn xiǎo} — My hands are small.
  \item \emph{Recall:} say the Mandarin for hot, then the Mandarin for to cry, then say \emph{shǒu} again
\end{itemize}

\begin{tcolorbox}[breakable,colback=teal!4,colframe=teal!35!black,title={Before you move on}]
What does \zh{手} mean? (\textquotedblleft{}Hand\textquotedblright{}.) Say it once more.
\end{tcolorbox}
\section[\texorpdfstring{\zh{手} (shǒu)}{手 (shǒu)}]{\zh{手} — a fall, two bars, a hook}

\label{lesson:ZH-W52-shou}



\begin{quote}
The word first, then the character inside it.

\begin{itemize}
\raggedright
  \item \emph{Recall:} say \emph{shǒu} — \textbf{R1}, one lesson back, before it decays
\end{itemize}

\begin{quote}
\zh{看} \zh{小}
\end{quote}

\emph{Read it:} each one aloud, with nothing said first — \emph{kàn}, \emph{xiǎo}
\end{quote}

\subsection*{Script — a fall, two bars, a hook}
\hlblockfigure{\detokenize{figures/ZH-W52-shou-filmstrip.pdf}}{How \zh{手} is written, stroke by stroke}

\begin{quote}
\zh{手}
\end{quote}

\textbf{One piece, four strokes, top to bottom.}

\textbf{A short left-falling stroke} across the top, from right to left.

\textbf{Two horizontals}, the lower one longer.

\textbf{Last, a vertical} down through both bars, hooking to the left at the foot.

The top of \zh{看} is this hand, squeezed flat over the eye.

\subsection*{Your turn}
\begin{itemize}
\raggedright
  \item \emph{Write it:} \zh{手}, all four, with the model in front of you
  \item \emph{Cover:} write it again from memory, and count the pieces, not the strokes
\end{itemize}

\begin{tcolorbox}[breakable,colback=teal!4,colframe=teal!35!black,title={Before you move on}]
How many strokes in \zh{手}? (\textbf{Four}.) Which way does the first stroke fall? (\textbf{To the left}.) Which character you know has a hand on top? (\textbf{\zh{看}}.)

Source: \href{https://raw.githubusercontent.com/chanind/hanzi-writer-data/68d10a4b21150cae5e1ebbd223eed289cf32d90c/data/\%E6\%89\%8B.json}{Hanzi Writer Data \zh{手}, pinned stroke paths and medians}.
\end{tcolorbox}
\section[tóu]{tóu — head}

\label{lesson:ZH-C52-tou}



\begin{quote}
Before the new one: say the Mandarin for to cry, then the Mandarin for hand.
\end{quote}

\subsection*{You'll want to know: tóu}
\emph{tóu} — \textquotedblleft{}head\textquotedblright{}.

\begin{quote}
\emph{tā de tóu hěn dà} — He's got a big head.
\end{quote}

Its characters come later; for now it is a word for the ear.

\begin{cousinweb}
\emph{tóu} is inherited from Middle Chinese \emph{duw} (Guangyun, in Baxter's transcription), glossed 'head'.
\end{cousinweb}

\begin{grammarlens}[title={What you've built}]
One more word for saying how you are.
\end{grammarlens}

\subsection*{Your turn}
\begin{itemize}
\raggedright
  \item \emph{Say it:} \emph{tóu}
  \item \emph{Say it:} \emph{tóu}, once more
  \item \emph{Say it:} \emph{tā de tóu hěn dà} — He's got a big head.
  \item \emph{Recall:} say the Mandarin for to cry, then the Mandarin for hand, then say \emph{tóu} again
\end{itemize}

\begin{tcolorbox}[breakable,colback=teal!4,colframe=teal!35!black,title={Before you move on}]
What does tóu mean? (\textquotedblleft{}Head\textquotedblright{}.) Say it once more.
\end{tcolorbox}
\section[jiǎo]{jiǎo — foot}

\label{lesson:ZH-C52-jiao}



\begin{quote}
Before the new one: say the Mandarin for hand, then the Mandarin for head.
\end{quote}

\subsection*{You'll want to know: jiǎo}
\emph{jiǎo} — \textquotedblleft{}foot\textquotedblright{}.

\begin{quote}
\emph{wǒ de jiǎo hěn lěng} — My feet are cold.
\end{quote}

Its characters come later; for now it is a word for the ear.

\begin{cousinweb}
\emph{jiǎo} is inherited from Middle Chinese \emph{kjak} (Guangyun, in Baxter's transcription), which the dictionary explains as the part that tucks back when you sit; whether that meant the leg or the foot is not settled.
\end{cousinweb}

\begin{grammarlens}[title={What you've built}]
One more word for saying how you are.
\end{grammarlens}

\subsection*{Your turn}
\begin{itemize}
\raggedright
  \item \emph{Say it:} \emph{jiǎo}
  \item \emph{Say it:} \emph{jiǎo}, once more
  \item \emph{Say it:} \emph{wǒ de jiǎo hěn lěng} — My feet are cold.
  \item \emph{Recall:} say the Mandarin for hand, then the Mandarin for head, then say \emph{jiǎo} again
\end{itemize}

\begin{tcolorbox}[breakable,colback=teal!4,colframe=teal!35!black,title={Before you move on}]
What does jiǎo mean? (\textquotedblleft{}Foot\textquotedblright{}.) Say it once more.
\end{tcolorbox}
\section[yǎnjing]{yǎnjing — eye}

\label{lesson:ZH-C52-yanjing}



\begin{quote}
Before the new one: say the Mandarin for head, then the Mandarin for foot.
\end{quote}

\subsection*{You'll want to know: yǎnjing}
\emph{yǎnjing} — \textquotedblleft{}eye\textquotedblright{}.

\begin{quote}
\emph{tā de yǎnjing hěn hǎokàn} — She has lovely eyes.
\end{quote}

Its characters come later; for now it is a word for the ear.

\begin{cousinweb}
\emph{yǎnjing} is built from yǎn 'eye' and jing 'eyeball', both inherited from Middle Chinese, \emph{ngeanX} and \emph{tsjeng} (Guangyun, in Baxter's transcription).
\end{cousinweb}

\begin{grammarlens}[title={What you've built}]
One more word for saying how you are.
\end{grammarlens}

\subsection*{Your turn}
\begin{itemize}
\raggedright
  \item \emph{Say it:} \emph{yǎnjing}
  \item \emph{Say it:} \emph{yǎnjing}, once more
  \item \emph{Say it:} \emph{tā de yǎnjing hěn hǎokàn} — She has lovely eyes.
  \item \emph{Recall:} say the Mandarin for head, then the Mandarin for foot, then say \emph{yǎnjing} again
\end{itemize}

\begin{tcolorbox}[breakable,colback=teal!4,colframe=teal!35!black,title={Before you move on}]
What does yǎnjing mean? (\textquotedblleft{}Eye\textquotedblright{}.) Say it once more.
\end{tcolorbox}
\section[ěrduo]{ěrduo — ear}

\label{lesson:ZH-C52-erduo}



\begin{quote}
Before the new one: say the Mandarin for foot, then the Mandarin for eye.
\end{quote}

\subsection*{You'll want to know: ěrduo}
\emph{ěrduo} — \textquotedblleft{}ear\textquotedblright{}.

\begin{quote}
\emph{háizi de ěrduo hěn xiǎo} — The child's ears are tiny.
\end{quote}

Its characters come later; for now it is a word for the ear.

\begin{cousinweb}
\emph{ěrduo} is built from ěr 'ear' and duo, both inherited from Middle Chinese, \emph{nyiX} and \emph{twaX} (Guangyun, in Baxter's transcription). The dictionary glosses the first as what hears and the second as something hanging from a tree; the second probably names the soft lobe that hangs, though that is not settled.
\end{cousinweb}

\begin{grammarlens}[title={What you've built}]
One more word for saying how you are.
\end{grammarlens}

\subsection*{Your turn}
\begin{itemize}
\raggedright
  \item \emph{Say it:} \emph{ěrduo}
  \item \emph{Say it:} \emph{ěrduo}, once more
  \item \emph{Say it:} \emph{háizi de ěrduo hěn xiǎo} — The child's ears are tiny.
  \item \emph{Recall:} say the Mandarin for foot, then the Mandarin for eye, then say \emph{ěrduo} again
\end{itemize}

\begin{tcolorbox}[breakable,colback=teal!4,colframe=teal!35!black,title={Before you move on}]
What does ěrduo mean? (\textquotedblleft{}Ear\textquotedblright{}.) Say it once more.
\end{tcolorbox}
\section[(dialogue)]{Looking back: from Telling the Time to Head, Hands and Feet}

\label{lesson:ZH-R52-review-48-52}



\begin{quote}
Say the Mandarin for o'clock and minute.
\end{quote}

\subsection*{Your turn}
\begin{itemize}
\raggedright
  \item \emph{Say it:} o'clock, minute, half, how many, to wait
  \item \emph{Say it:} before noon, noon, evening, afternoon, good evening
  \item \emph{Say it:} to sleep, to get up, to rest, health, to fall ill
  \item \emph{Say it:} happy, to be afraid, cold, hot, to cry
  \item \emph{Say it:} hand, head, foot, eye, ear
  \item \emph{Read it:} \zh{点} \zh{午} \zh{晚} \zh{手} — each one aloud, then a word that holds it
\end{itemize}

\begin{tcolorbox}[breakable,colback=teal!4,colframe=teal!35!black,title={Before you move on}]
O'clock? (\textbf{\zh{点}}, \emph{diǎn}.) Minute? (\emph{fēn}.) Half? (\emph{bàn}.) How many? (\emph{jǐ}.) To wait? (\emph{děng}.) Before noon? (\textbf{\zh{上午}}, \emph{shàngwǔ}.) Noon? (\textbf{\zh{中午}}, \emph{zhōngwǔ}.) Evening? (\textbf{\zh{晚上}}, \emph{wǎnshang}.) Afternoon? (\emph{xiàwǔ}.) Good evening? (\textbf{\zh{晚上好}}, \emph{wǎnshang hǎo}.) To sleep? (\emph{shuìjiào}.) To get up? (\emph{qǐchuáng}.) To rest? (\emph{xiūxi}.) Health? (\emph{shēntǐ}.) To fall ill? (\emph{bìng}.) Happy? (\emph{gāoxìng}.) To be afraid? (\emph{pà}.) Cold? (\emph{lěng}.) Hot? (\emph{rè}.) To cry? (\emph{kū}.) Hand? (\textbf{\zh{手}}, \emph{shǒu}.) Head? (\emph{tóu}.) Foot? (\emph{jiǎo}.) Eye? (\emph{yǎnjing}.) Ear? (\emph{ěrduo}.)
\end{tcolorbox}
